% F083｜定义域与正倒数的两类误区直接配图。
\clearpage
\begin{trapbox}
\textbf{\textcolor{CTrap}{易错一：以为“有指数或对数”就必定有奇偶性}}

只有特定结构才成立。例如 $a^x+1$ 通常\textbf{既不是奇函数也不是偶函数}；$\log_a x$ 的定义域仅为 $x>0$，也不能判为奇或偶。\textbf{必须先看结构，再核对定义。}

\textbf{\textcolor{CTrap}{易错二：不检查对数真数，直接用变号规律}}

$\ln\dfrac{kx-b}{kx+b}$ 不是对所有实数都有定义。当 $k,b\ne0$，只有 $|x|>|b/k|$ 才合法。\textbf{对数真数必须大于 $0$，不等于只要求分母不为 $0$。}

\end{trapbox}
\clearpage
\begin{trapbox}
\textbf{\textcolor{CTrap}{易错三：把 $\log_a|x|$ 与 $\log_a x$ 混淆}}
\begin{center}\begin{tikzpicture}[font=\small,>=stealth]
\draw[gray!40,rounded corners=3pt] (0,0) rectangle (14.2,8.4);
\node at (7.1,7.75) {绝对值改变了定义域：两幅图用相同坐标尺度};
\node at (3.65,6.8) {$y=\log_2x$};
\begin{scope}[shift={(3.65,3.85)},x=.82cm,y=.75cm]
\fill[gray!12] (-3,-2.6) rectangle (0,2.6);
\draw[->,gray] (-3,0)--(3,0) node[right] {$x$};
\draw[->,gray] (0,-2.6)--(0,2.6) node[above] {$y$};
\draw[blue,thick] plot coordinates {(0.1700000,-2.5563933) (0.1876875,-2.4135955) (0.2053750,-2.2836675) (0.2230625,-2.1644801) (0.2407500,-2.0543923) (0.2584375,-1.9521127) (0.2761250,-1.8566066) (0.2938125,-1.7670323) (0.3115000,-1.6826959) (0.3291875,-1.6030185) (0.3468750,-1.5275122) (0.3645625,-1.4557619) (0.3822500,-1.3874116) (0.3999375,-1.3221535) (0.4176250,-1.2597200) (0.4353125,-1.1998766) (0.4530000,-1.1424170) (0.4706875,-1.0871586) (0.4883750,-1.0339387) (0.5060625,-0.9826125) (0.5237500,-0.9330498) (0.5414375,-0.8851333) (0.5591250,-0.8387572) (0.5768125,-0.7938257) (0.5945000,-0.7502513) (0.6121875,-0.7079545) (0.6298750,-0.6668625) (0.6475625,-0.6269087) (0.6652500,-0.5880315) (0.6829375,-0.5501745) (0.7006250,-0.5132856) (0.7183125,-0.4773165) (0.7360000,-0.4422223) (0.7536875,-0.4079616) (0.7713750,-0.3744957) (0.7890625,-0.3417885) (0.8067500,-0.3098064) (0.8244375,-0.2785180) (0.8421250,-0.2478937) (0.8598125,-0.2179060) (0.8775000,-0.1885290) (0.8951875,-0.1597382) (0.9128750,-0.1315108) (0.9305625,-0.1038250) (0.9482500,-0.0766606) (0.9659375,-0.0499983) (0.9836250,-0.0238197) (1.0013125,0.0018923) (1.0190000,0.0271541) (1.0366875,0.0519811) (1.0543750,0.0763881) (1.0720625,0.1003890) (1.0897500,0.1239972) (1.1074375,0.1472253) (1.1251250,0.1700853) (1.1428125,0.1925887) (1.1605000,0.2147465) (1.1781875,0.2365692) (1.1958750,0.2580666) (1.2135625,0.2792484) (1.2312500,0.3001237) (1.2489375,0.3207013) (1.2666250,0.3409895) (1.2843125,0.3609963) (1.3020000,0.3807294) (1.3196875,0.4001963) (1.3373750,0.4194041) (1.3550625,0.4383594) (1.3727500,0.4570689) (1.3904375,0.4755389) (1.4081250,0.4937754) (1.4258125,0.5117843) (1.4435000,0.5295711) (1.4611875,0.5471413) (1.4788750,0.5645001) (1.4965625,0.5816525) (1.5142500,0.5986034) (1.5319375,0.6153574) (1.5496250,0.6319191) (1.5673125,0.6482929) (1.5850000,0.6644828) (1.6026875,0.6804931) (1.6203750,0.6963277) (1.6380625,0.7119904) (1.6557500,0.7274849) (1.6734375,0.7428147) (1.6911250,0.7579833) (1.7088125,0.7729941) (1.7265000,0.7878503) (1.7441875,0.8025551) (1.7618750,0.8171116) (1.7795625,0.8315226) (1.7972500,0.8457911) (1.8149375,0.8599199) (1.8326250,0.8739116) (1.8503125,0.8877689) (1.8680000,0.9014945) (1.8856875,0.9150906) (1.9033750,0.9285598) (1.9210625,0.9419045) (1.9387500,0.9551268) (1.9564375,0.9682290) (1.9741250,0.9812133) (1.9918125,0.9940818) (2.0095000,1.0068366) (2.0271875,1.0194795) (2.0448750,1.0320127) (2.0625625,1.0444378) (2.0802500,1.0567569) (2.0979375,1.0689717) (2.1156250,1.0810839) (2.1333125,1.0930953) (2.1510000,1.1050075) (2.1686875,1.1168222) (2.1863750,1.1285409) (2.2040625,1.1401651) (2.2217500,1.1516965) (2.2394375,1.1631364) (2.2571250,1.1744863) (2.2748125,1.1857476) (2.2925000,1.1969217) (2.3101875,1.2080099) (2.3278750,1.2190136) (2.3455625,1.2299339) (2.3632500,1.2407723) (2.3809375,1.2515298) (2.3986250,1.2622076) (2.4163125,1.2728070) (2.4340000,1.2833292) (2.4516875,1.2937751) (2.4693750,1.3041459) (2.4870625,1.3144428) (2.5047500,1.3246666) (2.5224375,1.3348185) (2.5401250,1.3448995) (2.5578125,1.3549105) (2.5755000,1.3648525) (2.5931875,1.3747265) (2.6108750,1.3845334) (2.6285625,1.3942740) (2.6462500,1.4039494) (2.6639375,1.4135602) (2.6816250,1.4231075) (2.6993125,1.4325920) (2.7170000,1.4420146) (2.7346875,1.4513760) (2.7523750,1.4606770) (2.7700625,1.4699185) (2.7877500,1.4791012) (2.8054375,1.4882258) (2.8231250,1.4972930) (2.8408125,1.5063036) (2.8585000,1.5152583) (2.8761875,1.5241577) (2.8938750,1.5330026) (2.9115625,1.5417936) (2.9292500,1.5505313) (2.9469375,1.5592165) (2.9646250,1.5678496) (2.9823125,1.5764314) (3.0000000,1.5849625)};
\end{scope}
\node at (3.65,1.25) {$D=(0,+\infty)$，不成对};
\node at (10.65,6.8) {$y=\log_2|x|$};
\begin{scope}[shift={(10.65,3.85)},x=.82cm,y=.75cm]
\draw[->,gray] (-3,0)--(3,0) node[right] {$x$};
\draw[->,gray] (0,-2.6)--(0,2.6) node[above] {$y$};
\draw[blue,thick] plot coordinates {(-3.0000000,1.5849625) (-2.9823125,1.5764314) (-2.9646250,1.5678496) (-2.9469375,1.5592165) (-2.9292500,1.5505313) (-2.9115625,1.5417936) (-2.8938750,1.5330026) (-2.8761875,1.5241577) (-2.8585000,1.5152583) (-2.8408125,1.5063036) (-2.8231250,1.4972930) (-2.8054375,1.4882258) (-2.7877500,1.4791012) (-2.7700625,1.4699185) (-2.7523750,1.4606770) (-2.7346875,1.4513760) (-2.7170000,1.4420146) (-2.6993125,1.4325920) (-2.6816250,1.4231075) (-2.6639375,1.4135602) (-2.6462500,1.4039494) (-2.6285625,1.3942740) (-2.6108750,1.3845334) (-2.5931875,1.3747265) (-2.5755000,1.3648525) (-2.5578125,1.3549105) (-2.5401250,1.3448995) (-2.5224375,1.3348185) (-2.5047500,1.3246666) (-2.4870625,1.3144428) (-2.4693750,1.3041459) (-2.4516875,1.2937751) (-2.4340000,1.2833292) (-2.4163125,1.2728070) (-2.3986250,1.2622076) (-2.3809375,1.2515298) (-2.3632500,1.2407723) (-2.3455625,1.2299339) (-2.3278750,1.2190136) (-2.3101875,1.2080099) (-2.2925000,1.1969217) (-2.2748125,1.1857476) (-2.2571250,1.1744863) (-2.2394375,1.1631364) (-2.2217500,1.1516965) (-2.2040625,1.1401651) (-2.1863750,1.1285409) (-2.1686875,1.1168222) (-2.1510000,1.1050075) (-2.1333125,1.0930953) (-2.1156250,1.0810839) (-2.0979375,1.0689717) (-2.0802500,1.0567569) (-2.0625625,1.0444378) (-2.0448750,1.0320127) (-2.0271875,1.0194795) (-2.0095000,1.0068366) (-1.9918125,0.9940818) (-1.9741250,0.9812133) (-1.9564375,0.9682290) (-1.9387500,0.9551268) (-1.9210625,0.9419045) (-1.9033750,0.9285598) (-1.8856875,0.9150906) (-1.8680000,0.9014945) (-1.8503125,0.8877689) (-1.8326250,0.8739116) (-1.8149375,0.8599199) (-1.7972500,0.8457911) (-1.7795625,0.8315226) (-1.7618750,0.8171116) (-1.7441875,0.8025551) (-1.7265000,0.7878503) (-1.7088125,0.7729941) (-1.6911250,0.7579833) (-1.6734375,0.7428147) (-1.6557500,0.7274849) (-1.6380625,0.7119904) (-1.6203750,0.6963277) (-1.6026875,0.6804931) (-1.5850000,0.6644828) (-1.5673125,0.6482929) (-1.5496250,0.6319191) (-1.5319375,0.6153574) (-1.5142500,0.5986034) (-1.4965625,0.5816525) (-1.4788750,0.5645001) (-1.4611875,0.5471413) (-1.4435000,0.5295711) (-1.4258125,0.5117843) (-1.4081250,0.4937754) (-1.3904375,0.4755389) (-1.3727500,0.4570689) (-1.3550625,0.4383594) (-1.3373750,0.4194041) (-1.3196875,0.4001963) (-1.3020000,0.3807294) (-1.2843125,0.3609963) (-1.2666250,0.3409895) (-1.2489375,0.3207013) (-1.2312500,0.3001237) (-1.2135625,0.2792484) (-1.1958750,0.2580666) (-1.1781875,0.2365692) (-1.1605000,0.2147465) (-1.1428125,0.1925887) (-1.1251250,0.1700853) (-1.1074375,0.1472253) (-1.0897500,0.1239972) (-1.0720625,0.1003890) (-1.0543750,0.0763881) (-1.0366875,0.0519811) (-1.0190000,0.0271541) (-1.0013125,0.0018923) (-0.9836250,-0.0238197) (-0.9659375,-0.0499983) (-0.9482500,-0.0766606) (-0.9305625,-0.1038250) (-0.9128750,-0.1315108) (-0.8951875,-0.1597382) (-0.8775000,-0.1885290) (-0.8598125,-0.2179060) (-0.8421250,-0.2478937) (-0.8244375,-0.2785180) (-0.8067500,-0.3098064) (-0.7890625,-0.3417885) (-0.7713750,-0.3744957) (-0.7536875,-0.4079616) (-0.7360000,-0.4422223) (-0.7183125,-0.4773165) (-0.7006250,-0.5132856) (-0.6829375,-0.5501745) (-0.6652500,-0.5880315) (-0.6475625,-0.6269087) (-0.6298750,-0.6668625) (-0.6121875,-0.7079545) (-0.5945000,-0.7502513) (-0.5768125,-0.7938257) (-0.5591250,-0.8387572) (-0.5414375,-0.8851333) (-0.5237500,-0.9330498) (-0.5060625,-0.9826125) (-0.4883750,-1.0339387) (-0.4706875,-1.0871586) (-0.4530000,-1.1424170) (-0.4353125,-1.1998766) (-0.4176250,-1.2597200) (-0.3999375,-1.3221535) (-0.3822500,-1.3874116) (-0.3645625,-1.4557619) (-0.3468750,-1.5275122) (-0.3291875,-1.6030185) (-0.3115000,-1.6826959) (-0.2938125,-1.7670323) (-0.2761250,-1.8566066) (-0.2584375,-1.9521127) (-0.2407500,-2.0543923) (-0.2230625,-2.1644801) (-0.2053750,-2.2836675) (-0.1876875,-2.4135955) (-0.1700000,-2.5563933)};
\draw[blue,thick] plot coordinates {(0.1700000,-2.5563933) (0.1876875,-2.4135955) (0.2053750,-2.2836675) (0.2230625,-2.1644801) (0.2407500,-2.0543923) (0.2584375,-1.9521127) (0.2761250,-1.8566066) (0.2938125,-1.7670323) (0.3115000,-1.6826959) (0.3291875,-1.6030185) (0.3468750,-1.5275122) (0.3645625,-1.4557619) (0.3822500,-1.3874116) (0.3999375,-1.3221535) (0.4176250,-1.2597200) (0.4353125,-1.1998766) (0.4530000,-1.1424170) (0.4706875,-1.0871586) (0.4883750,-1.0339387) (0.5060625,-0.9826125) (0.5237500,-0.9330498) (0.5414375,-0.8851333) (0.5591250,-0.8387572) (0.5768125,-0.7938257) (0.5945000,-0.7502513) (0.6121875,-0.7079545) (0.6298750,-0.6668625) (0.6475625,-0.6269087) (0.6652500,-0.5880315) (0.6829375,-0.5501745) (0.7006250,-0.5132856) (0.7183125,-0.4773165) (0.7360000,-0.4422223) (0.7536875,-0.4079616) (0.7713750,-0.3744957) (0.7890625,-0.3417885) (0.8067500,-0.3098064) (0.8244375,-0.2785180) (0.8421250,-0.2478937) (0.8598125,-0.2179060) (0.8775000,-0.1885290) (0.8951875,-0.1597382) (0.9128750,-0.1315108) (0.9305625,-0.1038250) (0.9482500,-0.0766606) (0.9659375,-0.0499983) (0.9836250,-0.0238197) (1.0013125,0.0018923) (1.0190000,0.0271541) (1.0366875,0.0519811) (1.0543750,0.0763881) (1.0720625,0.1003890) (1.0897500,0.1239972) (1.1074375,0.1472253) (1.1251250,0.1700853) (1.1428125,0.1925887) (1.1605000,0.2147465) (1.1781875,0.2365692) (1.1958750,0.2580666) (1.2135625,0.2792484) (1.2312500,0.3001237) (1.2489375,0.3207013) (1.2666250,0.3409895) (1.2843125,0.3609963) (1.3020000,0.3807294) (1.3196875,0.4001963) (1.3373750,0.4194041) (1.3550625,0.4383594) (1.3727500,0.4570689) (1.3904375,0.4755389) (1.4081250,0.4937754) (1.4258125,0.5117843) (1.4435000,0.5295711) (1.4611875,0.5471413) (1.4788750,0.5645001) (1.4965625,0.5816525) (1.5142500,0.5986034) (1.5319375,0.6153574) (1.5496250,0.6319191) (1.5673125,0.6482929) (1.5850000,0.6644828) (1.6026875,0.6804931) (1.6203750,0.6963277) (1.6380625,0.7119904) (1.6557500,0.7274849) (1.6734375,0.7428147) (1.6911250,0.7579833) (1.7088125,0.7729941) (1.7265000,0.7878503) (1.7441875,0.8025551) (1.7618750,0.8171116) (1.7795625,0.8315226) (1.7972500,0.8457911) (1.8149375,0.8599199) (1.8326250,0.8739116) (1.8503125,0.8877689) (1.8680000,0.9014945) (1.8856875,0.9150906) (1.9033750,0.9285598) (1.9210625,0.9419045) (1.9387500,0.9551268) (1.9564375,0.9682290) (1.9741250,0.9812133) (1.9918125,0.9940818) (2.0095000,1.0068366) (2.0271875,1.0194795) (2.0448750,1.0320127) (2.0625625,1.0444378) (2.0802500,1.0567569) (2.0979375,1.0689717) (2.1156250,1.0810839) (2.1333125,1.0930953) (2.1510000,1.1050075) (2.1686875,1.1168222) (2.1863750,1.1285409) (2.2040625,1.1401651) (2.2217500,1.1516965) (2.2394375,1.1631364) (2.2571250,1.1744863) (2.2748125,1.1857476) (2.2925000,1.1969217) (2.3101875,1.2080099) (2.3278750,1.2190136) (2.3455625,1.2299339) (2.3632500,1.2407723) (2.3809375,1.2515298) (2.3986250,1.2622076) (2.4163125,1.2728070) (2.4340000,1.2833292) (2.4516875,1.2937751) (2.4693750,1.3041459) (2.4870625,1.3144428) (2.5047500,1.3246666) (2.5224375,1.3348185) (2.5401250,1.3448995) (2.5578125,1.3549105) (2.5755000,1.3648525) (2.5931875,1.3747265) (2.6108750,1.3845334) (2.6285625,1.3942740) (2.6462500,1.4039494) (2.6639375,1.4135602) (2.6816250,1.4231075) (2.6993125,1.4325920) (2.7170000,1.4420146) (2.7346875,1.4513760) (2.7523750,1.4606770) (2.7700625,1.4699185) (2.7877500,1.4791012) (2.8054375,1.4882258) (2.8231250,1.4972930) (2.8408125,1.5063036) (2.8585000,1.5152583) (2.8761875,1.5241577) (2.8938750,1.5330026) (2.9115625,1.5417936) (2.9292500,1.5505313) (2.9469375,1.5592165) (2.9646250,1.5678496) (2.9823125,1.5764314) (3.0000000,1.5849625)};
\end{scope}
\node at (10.65,1.25) {$D=\mathbb R\setminus\{0\}$，成对};
\node at (7.1,.55) {左图没有负数输入；右图把正半图复制到左侧，但两图都不能代入 $0$。};
\end{tikzpicture}\end{center}


$\log_a|x|$ 的定义域为 $x\ne0$，左右对称，是偶函数；$\log_a x$ 只允许正数，既非奇也非偶。\textbf{绝对值不是装饰，而是改变了定义域与对称性。}
\end{trapbox}
\clearpage
\begin{trapbox}
\textbf{\textcolor{CTrap}{易错四：奇函数都必须经过原点？}}

只有在 $0$ 允许代入时，才有“奇函数在原点的函数值为 $0$”。例如
\[
 \frac{a^x+1}{a^x-1}
\]
在 $x=0$ 时分母为 $0$，它\textbf{仍是奇函数}，但不能说它经过 $(0,0)$。同样，两种 $\ln$ 比值函数在 $k,b\ne0$ 时也不允许 $x=0$。

\textbf{\textcolor{CTrap}{易错五：擅自给隐藏偶函数补上 $x=0$}}

\[
 x\left(\frac12+\frac1{2^x-1}\right)
\]
在 $x=0$ 时含有分母 $2^0-1=0$，原式无定义。\textbf{不能因前面乘了一个 $x$，就把它认作 $0$。}计算乘法前，各因子必须有定义。
\end{trapbox}
\clearpage
\begin{trapbox}
\textbf{\textcolor{CTrap}{易错六：忘记 $a^x$ 变成 $a^{-x}$，不是 $-a^x$}}

\[
 \boxed{a^{-x}=\frac1{a^x}},\quad\text{而不是}\quad -a^x.
\]
把这个最基本的变化记错，会把奇偶性判断全部颠倒。用 $2^{-1}=\tfrac12$ 与 $-2^1=-2$ 就能立即发现错误。

\textbf{\textcolor{CTrap}{易错七：只试 $x=1,-1$ 就当作完整证明}}

两个数的结果相等或相反，仅能帮助猜测，\textbf{不能替代对所有允许的 $x$ 证明}。正式证明用 $f(-x)$ 的通式，或使用已证明的奇偶函数运算规律。

\end{trapbox}
\clearpage
\begin{trapbox}
\textbf{\textcolor{CTrap}{易错八：把两个根号直接认为相反数}}
\begin{center}\begin{tikzpicture}[font=\small,>=stealth]
\draw[gray!40,rounded corners=3pt] (0,0) rectangle (14.2,8.4);
\node at (7.1,7.75) {两个根式真数都为正；互为倒数，取对数后才互为相反数};
\begin{scope}[shift={(4.4,1.8)},x=1.2cm,y=.75cm]
\draw[->,gray] (-2.8,0)--(2.8,0) node[right] {$x$};
\draw[->,gray] (0,0)--(0,6) node[above] {$y$};
\draw[gray,dotted] (-2.8,1)--(2.8,1);
\draw[blue,thick,solid] plot coordinates {(-2.7000000,0.1792360) (-2.6775000,0.1806473) (-2.6550000,0.1820804) (-2.6325000,0.1835356) (-2.6100000,0.1850134) (-2.5875000,0.1865145) (-2.5650000,0.1880392) (-2.5425000,0.1895883) (-2.5200000,0.1911621) (-2.4975000,0.1927614) (-2.4750000,0.1943866) (-2.4525000,0.1960385) (-2.4300000,0.1977176) (-2.4075000,0.1994247) (-2.3850000,0.2011603) (-2.3625000,0.2029252) (-2.3400000,0.2047200) (-2.3175000,0.2065456) (-2.2950000,0.2084027) (-2.2725000,0.2102920) (-2.2500000,0.2122145) (-2.2275000,0.2141708) (-2.2050000,0.2161619) (-2.1825000,0.2181887) (-2.1600000,0.2202521) (-2.1375000,0.2223530) (-2.1150000,0.2244925) (-2.0925000,0.2266715) (-2.0700000,0.2288910) (-2.0475000,0.2311523) (-2.0250000,0.2334563) (-2.0025000,0.2358043) (-1.9800000,0.2381975) (-1.9575000,0.2406370) (-1.9350000,0.2431242) (-1.9125000,0.2456604) (-1.8900000,0.2482469) (-1.8675000,0.2508853) (-1.8450000,0.2535769) (-1.8225000,0.2563233) (-1.8000000,0.2591260) (-1.7775000,0.2619868) (-1.7550000,0.2649072) (-1.7325000,0.2678890) (-1.7100000,0.2709341) (-1.6875000,0.2740444) (-1.6650000,0.2772217) (-1.6425000,0.2804681) (-1.6200000,0.2837857) (-1.5975000,0.2871767) (-1.5750000,0.2906433) (-1.5525000,0.2941879) (-1.5300000,0.2978129) (-1.5075000,0.3015208) (-1.4850000,0.3053142) (-1.4625000,0.3091959) (-1.4400000,0.3131686) (-1.4175000,0.3172352) (-1.3950000,0.3213988) (-1.3725000,0.3256626) (-1.3500000,0.3300298) (-1.3275000,0.3345037) (-1.3050000,0.3390879) (-1.2825000,0.3437860) (-1.2600000,0.3486019) (-1.2375000,0.3535394) (-1.2150000,0.3586026) (-1.1925000,0.3637957) (-1.1700000,0.3691231) (-1.1475000,0.3745894) (-1.1250000,0.3801993) (-1.1025000,0.3859577) (-1.0800000,0.3918696) (-1.0575000,0.3979402) (-1.0350000,0.4041751) (-1.0125000,0.4105798) (-0.9900000,0.4171603) (-0.9675000,0.4239224) (-0.9450000,0.4308725) (-0.9225000,0.4380169) (-0.9000000,0.4453624) (-0.8775000,0.4529158) (-0.8550000,0.4606842) (-0.8325000,0.4686749) (-0.8100000,0.4768955) (-0.7875000,0.4853536) (-0.7650000,0.4940572) (-0.7425000,0.5030145) (-0.7200000,0.5122337) (-0.6975000,0.5217236) (-0.6750000,0.5314929) (-0.6525000,0.5415504) (-0.6300000,0.5519052) (-0.6075000,0.5625668) (-0.5850000,0.5735443) (-0.5625000,0.5848475) (-0.5400000,0.5964858) (-0.5175000,0.6084690) (-0.4950000,0.6208069) (-0.4725000,0.6335092) (-0.4500000,0.6465856) (-0.4275000,0.6600460) (-0.4050000,0.6738999) (-0.3825000,0.6881569) (-0.3600000,0.7028264) (-0.3375000,0.7179176) (-0.3150000,0.7334393) (-0.2925000,0.7494003) (-0.2700000,0.7658089) (-0.2475000,0.7826729) (-0.2250000,0.8000000) (-0.2025000,0.8177971) (-0.1800000,0.8360709) (-0.1575000,0.8548271) (-0.1350000,0.8740714) (-0.1125000,0.8938082) (-0.0900000,0.9140418) (-0.0675000,0.9347755) (-0.0450000,0.9560120) (-0.0225000,0.9777531) (0.0000000,1.0000000) (0.0225000,1.0227531) (0.0450000,1.0460120) (0.0675000,1.0697755) (0.0900000,1.0940418) (0.1125000,1.1188082) (0.1350000,1.1440714) (0.1575000,1.1698271) (0.1800000,1.1960709) (0.2025000,1.2227971) (0.2250000,1.2500000) (0.2475000,1.2776729) (0.2700000,1.3058089) (0.2925000,1.3344003) (0.3150000,1.3634393) (0.3375000,1.3929176) (0.3600000,1.4228264) (0.3825000,1.4531569) (0.4050000,1.4838999) (0.4275000,1.5150460) (0.4500000,1.5465856) (0.4725000,1.5785092) (0.4950000,1.6108069) (0.5175000,1.6434690) (0.5400000,1.6764858) (0.5625000,1.7098475) (0.5850000,1.7435443) (0.6075000,1.7775668) (0.6300000,1.8119052) (0.6525000,1.8465504) (0.6750000,1.8814929) (0.6975000,1.9167236) (0.7200000,1.9522337) (0.7425000,1.9880145) (0.7650000,2.0240572) (0.7875000,2.0603536) (0.8100000,2.0968955) (0.8325000,2.1336749) (0.8550000,2.1706842) (0.8775000,2.2079158) (0.9000000,2.2453624) (0.9225000,2.2830169) (0.9450000,2.3208725) (0.9675000,2.3589224) (0.9900000,2.3971603) (1.0125000,2.4355798) (1.0350000,2.4741751) (1.0575000,2.5129402) (1.0800000,2.5518696) (1.1025000,2.5909577) (1.1250000,2.6301993) (1.1475000,2.6695894) (1.1700000,2.7091231) (1.1925000,2.7487957) (1.2150000,2.7886026) (1.2375000,2.8285394) (1.2600000,2.8686019) (1.2825000,2.9087860) (1.3050000,2.9490879) (1.3275000,2.9895037) (1.3500000,3.0300298) (1.3725000,3.0706626) (1.3950000,3.1113988) (1.4175000,3.1522352) (1.4400000,3.1931686) (1.4625000,3.2341959) (1.4850000,3.2753142) (1.5075000,3.3165208) (1.5300000,3.3578129) (1.5525000,3.3991879) (1.5750000,3.4406433) (1.5975000,3.4821767) (1.6200000,3.5237857) (1.6425000,3.5654681) (1.6650000,3.6072217) (1.6875000,3.6490444) (1.7100000,3.6909341) (1.7325000,3.7328890) (1.7550000,3.7749072) (1.7775000,3.8169868) (1.8000000,3.8591260) (1.8225000,3.9013233) (1.8450000,3.9435769) (1.8675000,3.9858853) (1.8900000,4.0282469) (1.9125000,4.0706604) (1.9350000,4.1131242) (1.9575000,4.1556370) (1.9800000,4.1981975) (2.0025000,4.2408043) (2.0250000,4.2834563) (2.0475000,4.3261523) (2.0700000,4.3688910) (2.0925000,4.4116715) (2.1150000,4.4544925) (2.1375000,4.4973530) (2.1600000,4.5402521) (2.1825000,4.5831887) (2.2050000,4.6261619) (2.2275000,4.6691708) (2.2500000,4.7122145) (2.2725000,4.7552920) (2.2950000,4.7984027) (2.3175000,4.8415456) (2.3400000,4.8847200) (2.3625000,4.9279252) (2.3850000,4.9711603) (2.4075000,5.0144247) (2.4300000,5.0577176) (2.4525000,5.1010385) (2.4750000,5.1443866) (2.4975000,5.1877614) (2.5200000,5.2311621) (2.5425000,5.2745883) (2.5650000,5.3180392) (2.5875000,5.3615145) (2.6100000,5.4050134) (2.6325000,5.4485356) (2.6550000,5.4920804) (2.6775000,5.5356473) (2.7000000,5.5792360)};
\fill[blue] (1,2.4142136) circle(2.3pt);
\draw[orange!85!black,thick,dashed] plot coordinates {(-2.7000000,5.5792360) (-2.6775000,5.5356473) (-2.6550000,5.4920804) (-2.6325000,5.4485356) (-2.6100000,5.4050134) (-2.5875000,5.3615145) (-2.5650000,5.3180392) (-2.5425000,5.2745883) (-2.5200000,5.2311621) (-2.4975000,5.1877614) (-2.4750000,5.1443866) (-2.4525000,5.1010385) (-2.4300000,5.0577176) (-2.4075000,5.0144247) (-2.3850000,4.9711603) (-2.3625000,4.9279252) (-2.3400000,4.8847200) (-2.3175000,4.8415456) (-2.2950000,4.7984027) (-2.2725000,4.7552920) (-2.2500000,4.7122145) (-2.2275000,4.6691708) (-2.2050000,4.6261619) (-2.1825000,4.5831887) (-2.1600000,4.5402521) (-2.1375000,4.4973530) (-2.1150000,4.4544925) (-2.0925000,4.4116715) (-2.0700000,4.3688910) (-2.0475000,4.3261523) (-2.0250000,4.2834563) (-2.0025000,4.2408043) (-1.9800000,4.1981975) (-1.9575000,4.1556370) (-1.9350000,4.1131242) (-1.9125000,4.0706604) (-1.8900000,4.0282469) (-1.8675000,3.9858853) (-1.8450000,3.9435769) (-1.8225000,3.9013233) (-1.8000000,3.8591260) (-1.7775000,3.8169868) (-1.7550000,3.7749072) (-1.7325000,3.7328890) (-1.7100000,3.6909341) (-1.6875000,3.6490444) (-1.6650000,3.6072217) (-1.6425000,3.5654681) (-1.6200000,3.5237857) (-1.5975000,3.4821767) (-1.5750000,3.4406433) (-1.5525000,3.3991879) (-1.5300000,3.3578129) (-1.5075000,3.3165208) (-1.4850000,3.2753142) (-1.4625000,3.2341959) (-1.4400000,3.1931686) (-1.4175000,3.1522352) (-1.3950000,3.1113988) (-1.3725000,3.0706626) (-1.3500000,3.0300298) (-1.3275000,2.9895037) (-1.3050000,2.9490879) (-1.2825000,2.9087860) (-1.2600000,2.8686019) (-1.2375000,2.8285394) (-1.2150000,2.7886026) (-1.1925000,2.7487957) (-1.1700000,2.7091231) (-1.1475000,2.6695894) (-1.1250000,2.6301993) (-1.1025000,2.5909577) (-1.0800000,2.5518696) (-1.0575000,2.5129402) (-1.0350000,2.4741751) (-1.0125000,2.4355798) (-0.9900000,2.3971603) (-0.9675000,2.3589224) (-0.9450000,2.3208725) (-0.9225000,2.2830169) (-0.9000000,2.2453624) (-0.8775000,2.2079158) (-0.8550000,2.1706842) (-0.8325000,2.1336749) (-0.8100000,2.0968955) (-0.7875000,2.0603536) (-0.7650000,2.0240572) (-0.7425000,1.9880145) (-0.7200000,1.9522337) (-0.6975000,1.9167236) (-0.6750000,1.8814929) (-0.6525000,1.8465504) (-0.6300000,1.8119052) (-0.6075000,1.7775668) (-0.5850000,1.7435443) (-0.5625000,1.7098475) (-0.5400000,1.6764858) (-0.5175000,1.6434690) (-0.4950000,1.6108069) (-0.4725000,1.5785092) (-0.4500000,1.5465856) (-0.4275000,1.5150460) (-0.4050000,1.4838999) (-0.3825000,1.4531569) (-0.3600000,1.4228264) (-0.3375000,1.3929176) (-0.3150000,1.3634393) (-0.2925000,1.3344003) (-0.2700000,1.3058089) (-0.2475000,1.2776729) (-0.2250000,1.2500000) (-0.2025000,1.2227971) (-0.1800000,1.1960709) (-0.1575000,1.1698271) (-0.1350000,1.1440714) (-0.1125000,1.1188082) (-0.0900000,1.0940418) (-0.0675000,1.0697755) (-0.0450000,1.0460120) (-0.0225000,1.0227531) (0.0000000,1.0000000) (0.0225000,0.9777531) (0.0450000,0.9560120) (0.0675000,0.9347755) (0.0900000,0.9140418) (0.1125000,0.8938082) (0.1350000,0.8740714) (0.1575000,0.8548271) (0.1800000,0.8360709) (0.2025000,0.8177971) (0.2250000,0.8000000) (0.2475000,0.7826729) (0.2700000,0.7658089) (0.2925000,0.7494003) (0.3150000,0.7334393) (0.3375000,0.7179176) (0.3600000,0.7028264) (0.3825000,0.6881569) (0.4050000,0.6738999) (0.4275000,0.6600460) (0.4500000,0.6465856) (0.4725000,0.6335092) (0.4950000,0.6208069) (0.5175000,0.6084690) (0.5400000,0.5964858) (0.5625000,0.5848475) (0.5850000,0.5735443) (0.6075000,0.5625668) (0.6300000,0.5519052) (0.6525000,0.5415504) (0.6750000,0.5314929) (0.6975000,0.5217236) (0.7200000,0.5122337) (0.7425000,0.5030145) (0.7650000,0.4940572) (0.7875000,0.4853536) (0.8100000,0.4768955) (0.8325000,0.4686749) (0.8550000,0.4606842) (0.8775000,0.4529158) (0.9000000,0.4453624) (0.9225000,0.4380169) (0.9450000,0.4308725) (0.9675000,0.4239224) (0.9900000,0.4171603) (1.0125000,0.4105798) (1.0350000,0.4041751) (1.0575000,0.3979402) (1.0800000,0.3918696) (1.1025000,0.3859577) (1.1250000,0.3801993) (1.1475000,0.3745894) (1.1700000,0.3691231) (1.1925000,0.3637957) (1.2150000,0.3586026) (1.2375000,0.3535394) (1.2600000,0.3486019) (1.2825000,0.3437860) (1.3050000,0.3390879) (1.3275000,0.3345037) (1.3500000,0.3300298) (1.3725000,0.3256626) (1.3950000,0.3213988) (1.4175000,0.3172352) (1.4400000,0.3131686) (1.4625000,0.3091959) (1.4850000,0.3053142) (1.5075000,0.3015208) (1.5300000,0.2978129) (1.5525000,0.2941879) (1.5750000,0.2906433) (1.5975000,0.2871767) (1.6200000,0.2837857) (1.6425000,0.2804681) (1.6650000,0.2772217) (1.6875000,0.2740444) (1.7100000,0.2709341) (1.7325000,0.2678890) (1.7550000,0.2649072) (1.7775000,0.2619868) (1.8000000,0.2591260) (1.8225000,0.2563233) (1.8450000,0.2535769) (1.8675000,0.2508853) (1.8900000,0.2482469) (1.9125000,0.2456604) (1.9350000,0.2431242) (1.9575000,0.2406370) (1.9800000,0.2381975) (2.0025000,0.2358043) (2.0250000,0.2334563) (2.0475000,0.2311523) (2.0700000,0.2288910) (2.0925000,0.2266715) (2.1150000,0.2244925) (2.1375000,0.2223530) (2.1600000,0.2202521) (2.1825000,0.2181887) (2.2050000,0.2161619) (2.2275000,0.2141708) (2.2500000,0.2122145) (2.2725000,0.2102920) (2.2950000,0.2084027) (2.3175000,0.2065456) (2.3400000,0.2047200) (2.3625000,0.2029252) (2.3850000,0.2011603) (2.4075000,0.1994247) (2.4300000,0.1977176) (2.4525000,0.1960385) (2.4750000,0.1943866) (2.4975000,0.1927614) (2.5200000,0.1911621) (2.5425000,0.1895883) (2.5650000,0.1880392) (2.5875000,0.1865145) (2.6100000,0.1850134) (2.6325000,0.1835356) (2.6550000,0.1820804) (2.6775000,0.1806473) (2.7000000,0.1792360)};
\fill[orange!85!black] (1,0.4142136) circle(2.3pt);
\draw[violet,densely dotted] (1,0)--(1,2.41421356);
\node[anchor=north] at (1,-.1) {$1$};
\end{scope}
\node[anchor=west,align=left,text width=4.5cm] at (8.9,4.3) {蓝实线：\\[4pt]$U(x)=\sqrt{x^2+1}+x$\\[10pt]橙虚线：\\[4pt]$V(x)=\sqrt{x^2+1}-x$\\[14pt]$U(1)=\sqrt2+1$\\[10pt]$V(1)=\sqrt2-1$};
\node at (7.1,.65) {$U(x)>0$、$V(x)>0$，且 $U(-x)=V(x)=1/U(x)$。};
\end{tikzpicture}\end{center}


$\sqrt{x^2+1}+x$ 与 $\sqrt{x^2+1}-x$ 不是互为相反数，它们\textbf{相乘等于 $1$}。所以取同一底数的对数后才\textbf{互为相反数}：
\[
 \log_a(\sqrt{x^2+1}-x)=-\log_a(\sqrt{x^2+1}+x).
\]
\end{trapbox}
\clearpage
\begin{trapbox}
\textbf{\textcolor{CTrap}{易错九：看到奇函数乘奇函数就仍判奇}}

\textbf{奇函数 $\times$ 奇函数 $=$ 偶函数}，前提是各函数都在关于 $0$ 对称的共同定义域上。例如 $x$ 与 $\dfrac{2^x+1}{2^x-1}$ 都为奇函数，它们的乘积就是偶函数。

\textbf{\textcolor{CTrap}{易错十：漏掉参数与底数条件}}

本条统一采用 $a>0$ 且 $a\ne1$。两种含 $k,b$ 的对数比值本条按 $k\ne0,b\ne0$ 讲解；若题目允许参数为 $0$，必须\textbf{重新检查定义域以及式子是否仍有意义}，不能机械套用速判表。

\textbf{一条防错流程：}\textbf{先看定义域是否左右对称，再看换号后是倒数、不变还是相反数，最后判断奇偶。}
\end{trapbox}
